% ---------------------------------------------------------------------------
% DRAFT: "Falsifiable predictions" section for main.tex
% Intended insertion point: after \section{Limitations and open questions}
% (line ~612, after the closing \end{itemize}), before \section{Conclusion}.
%
% Requires one new bib entry (see bottom of this file) — which finding #3 of
% the pre-submission review (engage SelectiveNet) already requires anyway.
% ---------------------------------------------------------------------------

\section{Falsifiable predictions}\label{sec:predictions}

The results above are proved on constructed uncertainty classes; whether
their domain includes real routed systems is an empirical question, and we
prefer to state now, before any data, what would settle it. Each prediction
below is quantitative, measurable on an instrumented system at modest cost,
and named together with the outcome that would refute it. We intend to test
them under a frozen pre-registration --- predictions, acceptance bands, and
decision rules committed publicly before any run; the bands sketched here
are indicative, not yet frozen.

\begin{description}
\item[P1 --- the coupling identity, measured.] Construct an explicit
  observation defect: a specialist whose routing decision must factor
  through a coarsened view $\varphi$ (redacted or truncated input), and a
  family of safe/fatal ensemble pairs that $\varphi$ provably cannot
  distinguish. Lemma~1 is an identity, not an inequality, so it predicts a
  tight band: at every threshold, the within-defect detector's gain-capture
  rate on the safe member equals its fatal-miss rate on the fatal member,
  within sampling error --- while a detector outside $\varphi$ separates the
  same pair by a margin bounded below by its measured advantage.
  \emph{Refuted if} within-defect detectors systematically deviate from
  equality --- in particular, capture above miss on the confounded family.
  That outcome would mean real coarsenings leak side information the formal
  $\varphi$ does not capture: $\sigma_0$ is generically positive in practice
  and the lemma's binding regime is empirically empty. We flag this as a
  live possibility, not a formality.

\item[P2 --- the observation/capacity dichotomy.] Same task, two defects:
  (a) an observation defect --- coarsened input, unrestricted detector
  class; (b) a capacity defect --- full input, restricted detector class.
  Intervention: transductive access, i.e.\ unlabeled samples of the
  deployment distribution. Theorems~1 and~2 jointly predict a two-by-two
  with no free cells: measurable rescue in (b) wherever the minimax gap of
  Theorem~2 is positive, and no rescue in (a) at any sample size.
  \emph{Refuted if} transductive access materially rescues (a), or fails to
  rescue (b) where the gap is positive.

\item[P3 --- the shared-representation boundary.] Rejectors that share
  representations with the actor empirically beat independently trained
  ones on average selective risk \citep{geifman2019selectivenet}; our
  results predict where that reverses, and we pre-commit to both halves.
  On standard benchmarks the shared rejector should \emph{win} the average
  case --- consistent with the prior literature, and no embarrassment to
  it. On a constructed confusable-fatal slice --- safe/fatal pairs the
  shared encoder maps together --- the shared rejector's selective risk is
  priced by the coupling identity and cannot fall below it, while an
  independent rejector with access to the distinction can. \emph{Refuted
  if} a shared-representation rejector beats the coupling bound on the
  constructed slice.

\item[P4 --- the premium as a quantitative law.] Theorem~1$'$ prices the
  surviving premium by the residual separability $\sigma_0$;
  Theorem~A gives the outside detector's premium $\Pi^*$ in measured
  quantities ($\bar\varepsilon$, $g$, $L$, $p$, $c_d$). Estimate $\sigma_0$
  and $\bar\varepsilon$ on held-out data, predict the realized premium
  across a grid of $(g, L)$, and regress realized on predicted. The
  prediction is a slope near one, not merely a positive sign.
  \emph{Refuted if} the slope departs systematically from one in the regime
  where the model's stated assumptions (i.i.d.\ periods, one-shot
  adversary) hold.
\end{description}

\noindent Two of these can fail in a way that damages the central claim
(P1, P3); the other two bound its scope (P2, P4). Stating them before any
data is collected is, we think, part of the result: a theory of deliberate
deficiency should declare its own fatal set.

% ---------------------------------------------------------------------------
% New bib entry required (references.bib) — also satisfies review finding #3:
%
% @inproceedings{geifman2019selectivenet,
%   author    = {Yonatan Geifman and Ran El-Yaniv},
%   title     = {{SelectiveNet}: A Deep Neural Network with an Integrated
%                Reject Option},
%   booktitle = {Proceedings of the 36th International Conference on Machine
%                Learning (ICML)},
%   year      = {2019},
%   pages     = {2151--2159}
% }
% ---------------------------------------------------------------------------
